\documentclass[tikz,border=6pt]{standalone}
\usepackage{peabra}
\begin{document}
\begin{tikzpicture}[x=1cm,y=1cm,
  fb/.style={doboz=22mm, font=\scriptsize, minimum height=11mm, inner sep=2pt},
  g/.style={circle, draw=peSotet, fill=white, line width=0.8pt, font=\scriptsize\bfseries, minimum size=7mm, inner sep=0}]
  \node[fb] (f1) at (0,0) {1.1 motorház és fedél befogása};
  \node[fb] (f2) at (2.9,0) {1.2 fedél pozicionálása};
  \node[g] (and1) at (5.0,0) {ÉS};
  \node[fb] (f3) at (7.1,0.9) {2.1–2.3 csavarok meghúzása, nyomatékmérés};
  \node[fb] (f4) at (7.1,-0.9) {következő darab előkészítése};
  \node[g] (and2) at (9.2,0) {ÉS};
  \node[g] (or1) at (10.6,0) {VAGY};
  \node[fb] (f5) at (12.8,0.9) {5.1–5.2 kivétel a kiadóba};
  \node[fb] (f6) at (12.8,-0.9) {3.2 hibás darab a selejt-dobozba};
  \node[g] (or2) at (14.9,0) {VAGY};
  \node[fb] (f7) at (16.9,0) {3.1, 3.3 naplózás, adatátadás};
  \draw[nyil] (f1) -- (f2); \draw[nyil] (f2) -- (and1);
  \draw[nyil] (and1) |- (f3); \draw[nyil] (and1) |- (f4);
  \draw[nyil] (f3) -| (and2); \draw[nyil] (f4) -| (and2);
  \draw[nyil] (and2) -- (or1);
  \draw[nyil] (or1) |- node[pos=0.3, left, font=\scriptsize, align=right] {nyomaték\\rendben} (f5);
  \draw[nyil] (or1) |- node[pos=0.3, left, font=\scriptsize, align=right] {tűrésen\\kívül} (f6);
  \draw[nyil] (f5) -| (or2); \draw[nyil] (f6) -| (or2);
  \draw[nyil] (or2) -- (f7);
  \draw[nyil, draw=peKiemel, dash pattern=on 3pt off 2pt] (f7.south) -- ++(0,-1.9) -| node[pos=0.25, below, font=\scriptsize, text=peKiemel] {következő ciklus} (f1.south);
\end{tikzpicture}
\end{document}
